% Copyright (c) 2026 Twin Vector Labs LLC.
% All rights reserved.
\documentclass[11pt]{article}

\usepackage[margin=1in]{geometry}
\usepackage{amsmath,amssymb,amsthm}
\usepackage{mathtools}
\usepackage{enumitem}
\usepackage{booktabs}
\usepackage{array}
\usepackage[colorlinks=true,linkcolor=blue,urlcolor=blue]{hyperref}

\newtheorem{principle}{Principle}
\theoremstyle{definition}
\newtheorem{definition}{Definition}
\newtheorem{remark}{Remark}
\newtheorem{finding}{Finding}

\newcommand{\C}{\mathbb{C}}
\newcommand{\R}{\mathbb{R}}
\newcommand{\Z}{\mathbb{Z}}
\newcommand{\rep}[1]{\mathbf{#1}}
\newcommand{\triv}{\mathbf{1}}
\newcommand{\abs}[1]{\left\lvert #1 \right\rvert}
\newcommand{\norm}[1]{\left\lVert #1 \right\rVert}
\newcommand{\sigmacol}{\sigma}
\newcommand{\Pin}{P_{\mathrm{in}}}
\newcommand{\Pout}{P_{\mathrm{out}}}
\newcommand{\dd}{\mathrm{d}}
\newcommand{\code}[1]{\texttt{#1}}
\DeclareMathOperator{\tr}{tr}
\DeclareMathOperator{\sgn}{sgn}

\emergencystretch=4em

\title{\bfseries Exploring the Flavor and Charge Sectors:\\
an epic synthesis of the emergent proton's electroweak content}
\author{Tessera --- cobordism subsystem
(epic \#410: \#411--\#418, \#429, \#437)}
\date{}

\begin{document}
\maketitle

\begin{abstract}
This note synthesizes the \emph{Flavor and Charge Sectors} exploration track (epic
\#410) end-to-end, now that its foundational and sector tickets have merged. The track
asked whether the proton's electroweak content --- electric charge, the proton/neutron
(flavor) distinction, and the spinor/current structure --- could be found by
\emph{splitting the behavior of the existing color register} (spatio-temporally and
across form-degree) rather than by bolting on parallel hole registers. The answer is now
substantially in hand, and it is more interesting than the plan anticipated. The two
geometric foundations landed: orientation is sourced from the standard simplicial formula
(\#412, \code{TemporalOrientation} / \code{ChainComplex::endSignCovector}), and the
asymmetric prism interior was replaced by a \emph{symmetric apex interior} that is now the
default (\#413, \code{Spacetime::symmetricStackCells}) --- driving the color-singlet
overlap to $1$ and the intertwining residual to machine zero, with a side finding that the
prism's ``spontaneous photon'' was triangulation-assisted. The charge sector is built and
emergent: the field strength splits into $E$ (timelike-leg) and $B$ (spacelike)
plaquettes (\#417), and electric charge reads as a temporal-sector Gauss-law holonomy
$Q=\oint\!\star F$ (\#411, \code{EigenstateSynthesis::gaussLawCharge}) that vanishes on a
Riemannian relaxation and populates only when Lorentzian worldlines do. The spin/current
sector is the Dirac--K\"ahler factor $(\dd+\delta)$ of the Hodge operator, with
$(\dd+\delta)^2=L$, exact Clifford relations, $\gamma^5$ chirality, the conserved current
$j^0$, and the $16=4\times4$ taste multiplicity (\#415, \code{DiracKahler}). The binding
question was reopened bipartitely (\#416): an orientation-reversing twist projects the
antisymmetric $\bar{\rep3}$ diquark, and a \emph{connected} two-body path reaches the
singlet (reachability $1.00$ vs.\ $0.55$ welded) --- confinement is not a fundamental
obstruction inside one connected bulk. The flavor sector returned a clean \emph{no-go}
(\#414): flavor as a relabeling-invariant split of the $A_4$-symmetric color register is
obstructed, because flavor is symmetry-odd while the singlet is symmetry-even and the two
cannot share a register; the fix is a symmetry-\emph{breaking} isospin doublet, a new
construction. The dimensional spike (\#418) and the dimension-generic apex interior (\#429,
\code{symmetricStackCells(nApexSlices)} with a rigorous \code{dualComplexIsValid} for
$n\ge4$, validated as a $4$-manifold over a tetrahedral $S^3$) scope the $3+1$\,D home and
expose the emergent twist as a PT-symmetric structure --- correctly read as
\emph{localized to creation/annihilation vertices} (St\"uckelberg--Feynman), not a
per-slice alternation. Finally, the track's findings force an \textbf{architecture pivot}:
the proton is assembled from \emph{sequences of bipartite junctions} relaxed in
\emph{one global step} with color-emergent creation, and the tripartite $W_{ABC}$ is
dropped. Charge lives in the emergent temporal sector, color in the register: the
charge\,$\otimes$\,color thesis. All numbers below are measured through the production
C++ and its Python suite; the open work (named, not created) is collected at the end.
\end{abstract}

\tableofcontents

\section{What the track set out to do}
\label{sec:intro}

The cobordism subsystem carries the proton's \emph{color-confinement skeleton}: on a
symmetric trivalent junction a color-symmetric quark input transports to the net-color-zero
singlet $\langle(1,\omega,\omega^2)\rangle$, confinement ($\sigmacol\to0$) is a Stokes
conservation law, and the singlet's location is forced by the geometry's $\Z_3$ symmetry
rather than imposed. What it did \emph{not} carry was the electroweak content: charge,
the proton/neutron flavor distinction, and the spinor/current structure. The earlier
conclusion was a \emph{richer construction} --- a flavor register parallel to color, plus
spinor fields. Epic \#410 took the opposite road.

\begin{principle}[Split the register, do not multiply it]
\label{prin:split}
The proton's electroweak content is sought by splitting the behavior of the
\emph{existing} register --- across causal type (temporal vs.\ spatial) and across
form-degree (the inhomogeneous cochain) --- rather than by adding parallel hole registers.
Independence between sectors is supplied by orthogonality in the spectral/causal
decomposition, not by disjoint topology.
\end{principle}

This synthesis records what that program produced. The plan note (\#420,
\code{flavor\allowbreak\_charge\allowbreak\_sectors.tex}) laid out the sectors and their
falsifiable tests; the preliminary results note
(\code{flavor\allowbreak\_charge\allowbreak\_results.tex}) recorded the first wave. Here
we close the loop: every sector ticket has now landed, two of the load-bearing hypotheses
were \emph{killed} (the flavor split; the per-slice twist reading), the binding question
was answered, and the assembly architecture pivoted. Each section names the merged ticket,
the production API, and the measured verdict.

\section{Geometric foundations}
\label{sec:foundations}

The two foundations removed arbitrariness from the geometry; every downstream observable
inherits the symmetry they restored.

\subsection{Orientation is geometric, not a vertex-label sort (\#412)}
\label{sec:orient}

The color charge of a window is a sum of per-hole periods. The inputs were already read
in the surface's induced orientation, but the \emph{emergent result} was summed in each
hole's own sorted-vertex orientation --- a function of the arbitrary vertex numbering, so
a relabeling could flip a hole's term with no conjugating partner and scramble the sum.
The fix sources every sign from the standard simplicial orientation. Two distinct notions
of ``orientation'' were separated in code: the Causal-Dynamical-Triangulations $(t_i,t_f)$
temporal class became \code{TemporalOrientation} (with \code{orientationOf} the
standard-orientation anchor), reserving the \emph{spatial} induced orientation for
\code{ChainComplex::endSignCovector} (a globally coherent $\pm1$ covector propagated
across shared facets). Applying the result block's induced signs symmetrically with the
inputs turns the signed carried charge into the geometry's Stokes invariant:

\begin{center}
\small
\begin{tabular}{@{}lcc@{}}
\toprule
input & signed $\abs{\sigmacol_R}$ (induced) & raw $\abs{\sigmacol_R}$ (sorted) \\
\midrule
neutral ($q\bar q$ triple, $\Sigma=0$) & $0.0000$ (exact) & $0.2258$\\
colored ($R,G,B$, $\Sigma=1$ each)     & $3.0000$ (exact) & $1.0219$\\
\bottomrule
\end{tabular}
\end{center}

\noindent The signed sum is exactly $\sigmacol_R=-(\sigmacol_A+\sigmacol_B+\sigmacol_C)$;
the raw sum is neither, and is label-dependent. Correctness is pinned by a
\emph{relabeling-invariance} test (the track's $G_6$): under seeded permutations of the
base-surface vertex ids the signed $\sigmacol_R$, singlet overlap, and intertwining
residual are invariant up to one global sign shared by all windows, while the raw sum
drifts.

\subsection{The symmetric apex interior is the default; the singlet is exact (\#413)}
\label{sec:apex}

The trivalent junction was built by extruding the holed base surface $\times I$ and
tetrahedralizing each triangular prism by the Freudenthal staircase, whose diagonal is
chosen by a vertex-label sort. That choice --- not the metric --- was the documented
source of the intertwining residual $\norm{M\Pin-\Pout M}/\norm{M}\approx 4.26\times
10^{-2}$ that held the singlet at $0.999$. The replacement, \code{Spacetime::}\allowbreak
\code{symmetricStackCells}, cones each base triangle up to a single face-apex vertex (a
$(3,1)$ tetrahedron) and down to the top copy (a $(1,3)$ tetrahedron); the octahedral gap
over a shared base edge is split along the \emph{canonical dual edge} $f_1f_2$ (the two
face-centres), fixed by the incident faces rather than by ids. With no label-sort anywhere
the subdivision is equivariant under the icosahedral symmetry, and the transport
intertwines the color $\Z_3$ \emph{exactly}:

\begin{center}
\small
\begin{tabular}{@{}lcc@{}}
\toprule
 & prism (Freudenthal) & \textbf{symmetric apex (default)}\\
\midrule
intertwining residual $\norm{M\Pin-\Pout M}/\norm{M}$ & $4.26\times10^{-2}$ & $\mathbf{4.5\times10^{-14}}$\\
color-singlet overlap & $0.999$ & $\mathbf{1.000000}$\\
valid manifold (\code{dualComplexValid}) / $b_1$ & yes / $11$ & yes / $11$\\
\bottomrule
\end{tabular}
\end{center}

\noindent The symmetric interior is the default; \code{set\_symmetric\_interior(false)}
selects the legacy prism. The twenty existing junction tests pass unchanged --- the
residual simply moved inside their tolerances --- so the color-confinement skeleton
(charge conservation exact at $\sim10^{-14}$, $\sigmacol\to0$, the $\Z_3$ spectrum
$\{1,\omega,\omega^2\}$) is preserved while the singlet becomes exact.

\begin{finding}[The spontaneous photon was triangulation-assisted]
\label{find:photon}
On the prism, making the cross-layer worldlines timelike and relaxing drives one worldline
through null ($\ell^2\to0$) --- the predicted ``photon.'' On the symmetric interior the
same seed does \emph{not} produce it: through $120$ relaxation iterations all $450$
worldline edges stay uniformly timelike at $\ell^2=-0.3$ with \emph{zero} near-null edges.
The prism's photon was triggered by the triangulation asymmetry singling out one edge; the
symmetric relaxation preserves the worldline symmetry. A genuine photon should emerge from
a symmetry-breaking \emph{source}, not the bare uniform seed --- the symmetric interior is
the more honest setting, and a source-driven photon is the natural follow-up.
\end{finding}

\section{The charge sector}
\label{sec:charge}

\subsection{$E$ and $B$ from the temporal/spatial split (\#417)}
\label{sec:eb}

Let $A\in\Omega^1$ be a $U(1)$ connection and $F=\dd A\in\Omega^2$ its curvature. On the
Lorentzian complex each plaquette is classified by the causal type of its edges:
\begin{equation}
  F = \underbrace{F^{(E)}}_{\text{a timelike leg}}
      \;\oplus\;
      \underbrace{F^{(B)}}_{\text{purely spacelike}} ,
\end{equation}
the discrete analogue of $E_i=F_{0i}$, $B_{ij}=F_{ij}$. This is implemented as
\code{EigenstateSynthesis::}\allowbreak\code{curvatureFromConnection} ($F=\dd A$, the same induced-orientation
signed edge sum the periods use) and \code{fieldStrengthSplit}, returning the electric and
magnetic parts and their disjoint, complete cell index lists. The tests confirm the split
tracks the causal structure: on a purely spacelike (static) seed $\norm{E}=0$ exactly and
$B=F$ on every cell; under timelike flux $\norm{E}>0$ and the electric cells equal the
independently-recomputed timelike-leg plaquettes; $E+B=F$ to $10^{-12}$; and $F=\dd A$ is
gauge invariant ($A\to A+\dd\chi$ leaves $E,B$ unchanged). The reader is purely
observational --- the proton sector is bit-for-bit unchanged with it present or absent ---
and because the split is read off the cochain by causal type it is
\emph{interior-shape-independent}: it holds equally on the prism and the symmetric apex.

\subsection{Electric charge as a Gauss-law holonomy (\#411)}
\label{sec:gauss}

Electric charge is the discrete Gauss law, the flux of $E$ through a closed surface
enclosing the quark windows:
\begin{equation}
  Q \;=\; \oint_{\partial V} \star F \;=\; \int_V \star j ,
  \label{eq:gauss}
\end{equation}
implemented as \code{EigenstateSynthesis::gaussLawCharge}. The key property is what kind
of object this is. The previously-considered ``flavor-weighted covector''
$(\tfrac23,\tfrac23,-\tfrac13)$ is not $A_4$-equivariant and is not topologically
protected; it would degrade under a non-symmetric seed the way charge conservation
degrades under jitter. The Gauss-law holonomy \eqref{eq:gauss} is a genuine gauged-$U(1)$
invariant --- emergent and metric-robust.

\begin{finding}[Charge is emergent and Lorentzian, not hand-weighted]
\label{find:charge}
$Q$ is \emph{not} a hand-weighted covector. It is read off the relaxed geometry, and it is
\emph{zero on an all-spacelike (Riemannian) relaxation}: with no timelike worldlines the
$E$-sector is empty and Gauss's law returns no charge. The charge populates only when
Lorentzian worldlines do --- charge is intrinsically a temporal-sector object, exactly as
the reframe predicted. This is the \emph{charge}\,$\otimes$\,\emph{color} thesis in one
fact: color is the topological register kernel (metric-robust, always present); charge is
the temporal Gauss-law holonomy (present only with a Lorentzian causal structure).
\end{finding}

\noindent In the Standard Model the fractional part of the quark charge is tied to color
through the center $\Z_3$ (triality) via $Q=I_3+\tfrac12Y$. The construction already
carries triality as $\Pout$; whether the $\tfrac13$-quantum of the Gauss-law charge
\emph{is} that triality is the natural cross-check, alongside the agreement of $Q$ with the
Dirac--K\"ahler current $j^0$ of \S\ref{sec:dk}.

\section{The spin / current sector: Dirac--K\"ahler (\#415)}
\label{sec:dk}

The inhomogeneous cochain $\Omega^\bullet=\bigoplus_k\Omega^k$ is the discrete
Dirac--K\"ahler field, and $D=\dd+\delta$ is its Dirac operator (the class
\code{DiracKahler}). Assembled from the integer boundary maps and metric weights ---
reusing the convention the Hodge Laplacian already squares --- it reproduces the expected
algebra to machine precision:

\begin{center}
\small
\begin{tabular}{@{}lcl@{}}
\toprule
property & measured & meaning\\
\midrule
$(\dd+\delta)^2 = L$ & residual $\sim 6\times10^{-15}$ & $D$ is a square root of the Hodge Laplacian\\
$\{\gamma^a,\gamma^b\}=2\eta^{ab}$ & residual $0.0$ (exact) & the Clifford algebra of the $1$-cochain action\\
$\gamma^5$ chirality & $(\gamma^5)^2=I$, $\{\gamma^5,\gamma^a\}=0$ & the left/right (chiral) split\\
$j^0=\bar\Phi\gamma^0\Phi$ & $=\pm\,\Phi^\dagger\Phi$ & the conserved (Noether) charge density\\
multiplicity & $=4$ & the taste degeneracy $\Lambda(\C^4)=16=4\times4$\\
\bottomrule
\end{tabular}
\end{center}

\noindent The gamma generators are the $16\times16$ Clifford action of unit $1$-cochains
on the form fiber, $\gamma^a=\varepsilon(e^a)+\eta^{aa}\iota(e^a)$ (an edge raises and
lowers form degree) --- the literal realization of ``gammas $=$ the Clifford action of
edges.'' In four dimensions $\Omega^\bullet\cong\C l(1,3)\otimes\C\cong M_4(\C)$, so the
$4\times4$ Dirac structure is \emph{factored out of the operator the construction already
builds}. The charge density is \code{DiracKahler::chargeDensity}, the time component of
the conserved current; its cross-check against the Gauss-law charge of \S\ref{sec:gauss}
ties the operator sector to the temporal sector. The $16=4\times4$ taste multiplicity (the
lattice Dirac--K\"ahler ``doubling'') was the plan's candidate flavor index --- a
hypothesis the flavor sector then tested and rejected (\S\ref{sec:flavor}).

\section{The binding structure: a bipartite revisit (\#416)}
\label{sec:bipartite}

The proton was originally placed on a trivalent junction because the color singlet
$\varepsilon_{ijk}$ is absent from $\rep3\otimes\rep3=\rep6\oplus\bar{\rep3}$ and first
appears in $\rep3^{\otimes3}$. That theorem rules out a \emph{generic} bipartite merge,
but not two bipartite steps with a projection: if the first merge projects onto the
antisymmetric diquark $\bar{\rep3}=\varepsilon qq$, then
$\bar{\rep3}\otimes\rep3=\triv\oplus\rep8$ contains the singlet. The prism extrusion
already accepted a cumulative \emph{twist} permutation, an orientation-reversing tube that
is a geometric antisymmetrizer; \#416 exposed it cleanly as
\code{RegisterTopology::orientationReversingTwist} / \code{setTwist} with
\code{Cobordism::twistedCylinder}.

Two results landed. First, the twisted bipartite tube projects the emergent diquark onto
$\bar{\rep3}$ (the antisymmetric channel) rather than the generic symmetric average ---
the orientation-reversing twist \emph{is} the geometric antisymmetrizer the
representation theory wanted. Second, and more consequential:

\begin{finding}[Confinement is not fundamental in a connected bulk]
\label{find:reach}
The deeper obstruction was that the colored diquark intermediate is non-carriable
(confinement): only $\sigmacol=0$ combinations carry as \emph{free boundary} states. A
\emph{connected} two-body path geometry --- where the diquark is internal and never a free
boundary --- evades this. Measured, the connected path reaches the singlet with
reachability $1.00$, versus $0.55$ for the welded variant. The colored intermediate's
non-carriability is therefore \emph{not} a fundamental wall: it is hosted in the bulk of one
connected relaxation. This single fact reshapes the assembly architecture (\S\ref{sec:pivot}).
\end{finding}

\section{The flavor sector: a no-go (\#414)}
\label{sec:flavor}

Proton (uud) and neutron (udd) differ only by flavor; the color-only geometry cannot tell
them apart. Principle~\ref{prin:split} sought flavor as a split of the existing register.
That hope is now \emph{refuted}, cleanly and for a structural reason.

\begin{principle}[Flavor cannot share the color register]
\label{prin:flavor-nogo}
Any relabeling-invariant per-window measurable commutes with the window-cycling
automorphism $g$ that realizes the color $\Z_3$, hence is constant on $g$'s orbit. The
uud/udd discriminator is a function on that orbit, so it collapses: it cannot distinguish
windows that $g$ identifies. Measured, the discriminator margin is $\sim 3\times10^{-6}$,
bounded by the intertwining residual itself --- i.e.\ machine zero, no signal.
\end{principle}

The reason is symmetry parity. The color singlet is the symmetry-\emph{even} ($A_4$/$\Z_3$
invariant) line; flavor is symmetry-\emph{odd} (it must \emph{distinguish} the windows the
symmetry identifies). A single register cannot host both an even invariant and an odd
discriminator on the same orbit. Neither candidate mechanism survives: the spatio-temporal
split is relabeling-invariant (so it is constant on the orbit), and the Dirac--K\"ahler
taste multiplicity is an operator-sector degeneracy, not a per-window state label.

\begin{finding}[The fix is symmetry-breaking isospin, not a dimensional artifact]
\label{find:flavor}
The collapse is \emph{not} an artifact of working in $2+1$\,D or of the discretization ---
it is forced by the equivariance of the read-out. The cure is a symmetry-\emph{breaking}
isospin doublet: a genuinely new construction in which the $u$/$d$ asymmetry is put in as a
broken $SU(2)$ that \emph{acts} (the ladder $p\leftrightarrow n$), exactly as the color
$\Z_3$ acts through $\Pout$. A static even register cannot be coaxed into producing an odd
label; one must break the symmetry on purpose. This is the track's sharpest negative
result and it redirects the flavor follow-up away from ``find it in the existing register''
toward ``build the broken-isospin doublet.''
\end{finding}

\section{The dimensional question (\#418, \#429)}
\label{sec:dim}

Full $\vec E,\vec B$ and the full $4\times4$ Dirac algebra live only in $3+1$ dimensions ---
a genuine $S^3$ spatial slice extruded in time. The present junction is $2+1$\,D
($S^2\times I$), a reduced sector. \#418 is the spike that scopes the $3+1$\,D home and
returns a go/no-go for the full electroweak content; \#429 makes the apex interior
dimension-generic so the spike has a buildable interior.

\subsection{Dimension-generic apex interior and an $n\ge4$ manifold check (\#429)}

\code{Spacetime::symmetricStackCells} now carries an \code{nApexSlices} parameter
(default $1$, reproducing the single-reflection \#413 interior): each apex is a centre of
point-reflection, its neighbour edges inverted through it to form the next slice plus a
cap, stacked \code{nApexSlices} times. The interior connectivity generalizes by
\emph{coface-mirroring} --- cone each top $d$-simplex to an apex and reflect
apex-connectivity across a $(d-1)$-facet shared by two cofaces. The dimension-genericity
is backed by a rigorous \code{ChainComplex::dualComplexIsValid} that extends to $n\ge4$;
the construction is \emph{validated as a $4$-manifold} over a tetrahedral $S^3$ base.

\begin{finding}[The emergent twist is PT-symmetric, localized to creation/annihilation]
\label{find:twist}
Each apex is a $-1$ point-reflection, i.e.\ a PT inversion. Composed over the stack, the
reflect-and-cap iteration tracks the $\gamma^5$ chiral split (an $8/8$ partition of the
$16$ form components) and carries the net $\Z_2$ spinor sign ($2\pi\mapsto-1$): the
geometry realizes the spinor double cover, not merely the orientation $\Z_2$.
\emph{Corrected reading, important:} the naive per-slice alternation $(-1)^j$
\textbf{over-reads} the structure. For a \emph{freely propagating} worldline the apex cap
is manifold-closure only --- it carries no physical time reversal. The genuine PT /
time-orientation reversal is \emph{localized to creation/annihilation vertices}: an
antiparticle is a particle running backward in time (St\"uckelberg--Feynman), so the PT
flip belongs at the pair-creation/annihilation event, not at every slice of a propagating
line. Reading a twist on every slice would manufacture spurious sign structure along free
propagation; the physical content is concentrated at the events.
\end{finding}

\section{CI hygiene (\#437)}
\label{sec:ci}

Making the symmetric apex the default (\#413/\#429) and renaming the orientation class
(\#412) left five post-merge test failures. These were \emph{stale prism-era tests}, not
\#429 regressions: they asserted the prism's behavior in places where the default geometry
had moved underneath them. The golden invariants $G_1$--$G_5$ stayed green throughout. The
fixes pin the tests to the right path: the photon test now exercises the \emph{prism} path,
where the spontaneous photon is real (Finding~\ref{find:photon}); the residual test asserts
the machine-zero floor of the symmetric default; and the relabeling test pins the prism
stride. The track's shared regression module
(\code{tests/cobordism/test\_epic410\_invariants.py}) remains the gate every PR must keep
green.

\section{Architecture pivot: bipartite junctions and one global relaxation}
\label{sec:pivot}

The track's binding result (Finding~\ref{find:reach}) forces a change in how the proton is
\emph{assembled}. The tripartite $W_{ABC}$ (\code{TripartiteRegisterTopology}) is
\textbf{dropped} from the assembly path. The proton is built from \emph{sequences of
bipartite (two-body) junctions} relaxed in \emph{one global step}, with color-emergent
creation.

\paragraph{Why both at once.} The colored diquark intermediate is floored by confinement if
it is \emph{pinned} as a free boundary state (a staged, separate-relaxation assembly fails
exactly there). But it is \emph{hosted} in the bulk of one connected relaxation --- this is
Finding~\ref{find:reach}'s reachability $1.00$. So bipartite structure and
global-relaxation are required \emph{together}: pin only the initial and final states,
relax the entire interior at once, and let the diquark and other intermediates
\emph{emerge}. Staging is at most a warm start, never a separate relaxation.

\paragraph{Charge $\otimes$ color.} The pivot keeps the two sectors in their natural homes.
Charge stays in the emergent temporal sector (the Gauss-law holonomy of
\S\ref{sec:gauss}, Finding~\ref{find:charge}); color stays in the register kernel. The
assembly composes them without a parallel charge register --- charge is read emergently
from the relaxed Lorentzian geometry, not carried as a hand-weighted covector.

\paragraph{Creation, with the old assumptions quarantined.} The assembly reuses only the
\emph{creation event} of the existing \code{pairCreate} primitive. The Charged-Cartan
assumptions that surrounded it --- a parallel $\hat Q$ vertex register, an imposed
pair-Hamiltonian, Metropolis Monte-Carlo dynamics, a hand-dialed CP phase --- are
\textbf{quarantined}: they are not part of the emergent path. Creation is color-emergent
(color-indefinite at birth, crystallizing from context) and charge is read by Gauss's law.

\section{Faithfulness: the golden invariants held}
\label{sec:faithful}

Every result above landed under a fixed set of proton invariants, the track's
$G_1$--$G_8$, shared in one regression module so later subtasks could not silently regress
earlier ones.

\begin{center}
\small
\setlength{\tabcolsep}{5pt}
\begin{tabular}{@{}clp{0.62\linewidth}@{}}
\toprule
 & invariant & status across the epic\\
\midrule
$G_1$ & proton present & singlet overlap $\to 1.000000$ on the symmetric default (\#413)\\
$G_2$ & confinement & $\sigmacol\to0$ for the symmetric input\\
$G_3$ & charge conservation & $\sigmacol_R=-(\sigmacol_A+\sigmacol_B+\sigmacol_C)$ exact ($\sim10^{-14}$)\\
$G_4$ & topology & $b_1=11$ and a valid manifold; \code{dualComplexIsValid} now $n\ge4$ (\#429)\\
$G_5$ & color $\Z_3$ intact & $\Pout$ eigenvalues $\{1,\omega,\omega^2\}$\\
$G_6$ & relabeling-invariance & owned by \#412; the no-go of \#414 is a \emph{consequence} of it\\
$G_7$ & determinism & fixed seed $\Rightarrow$ reproducible; no wall-clock/random seeding\\
$G_8$ & emergent-first & charge (\#411), diquark (\#416), twist (\#429) all read off the relaxed geometry\\
\bottomrule
\end{tabular}
\end{center}

\noindent The geometry change (\#413) preserved the color sector exactly while sharpening
the singlet; the orientation change (\#412) preserved the manifold and spectrum while
making the charge a label-independent invariant; the Dirac--K\"ahler and $E/B$ additions
(\#415, \#417) are observable-only. No golden result drifted. Notably, $G_6$ is not merely
a guard here --- the flavor no-go (\S\ref{sec:flavor}) is \emph{derived} from it: the same
relabeling-invariance that makes the charge faithful is what forbids flavor from hiding in
the existing register.

\section{What is established, and what is open}
\label{sec:assess}

\paragraph{Established (merged and measured).} The geometry is honest (no label-sort
artifacts, \#412); the proton color singlet is \emph{exact} (overlap $1.000000$, residual
$4.5\times10^{-14}$, \#413); the field strength carries a faithful $E/B$ split (\#417);
electric charge is an emergent, metric-robust Gauss-law holonomy that vanishes without
Lorentzian worldlines (\#411); the Dirac--K\"ahler operator gives exact Clifford structure,
$\gamma^5$ chirality, a conserved current, and the $4\times4$ taste multiplicity (\#415);
a twisted bipartite tube projects the antisymmetric diquark and a connected path reaches
the singlet at reachability $1.00$ (\#416); flavor as a register split is \emph{ruled out}
with the cure identified (\#414); and the apex interior is dimension-generic, $4$-manifold
validated over $S^3$, with the emergent twist correctly localized to
creation/annihilation events (\#418, \#429). CI is green and pinned (\#437).

\paragraph{Open work (named here, not created).}
\begin{itemize}[nosep]
  \item \textbf{\#435 --- bipartite color-emergent creation primitive + charge\,$\leftrightarrow$\,color
  bridge.} The foundation of the assembly pivot: a two-body $q/\bar q$ creation event,
  color-indefinite at birth, with charge read by Gauss's law and the Charged-Cartan
  assumptions quarantined. Blocks \#434 and \#438.
  \item \textbf{\#434 --- Experiment A: emergent intermediates.} The whole
  $q/\bar q\to$ proton \& anti-proton event as one connected bipartite cobordism; pin only
  initial and final, relax the entire interior at once, let the diquark and intermediates
  emerge; watch for sourced photons.
  \item \textbf{\#438 --- Experiment B: fixed bipartite sequence + A-vs-B comparison.}
  Reuse \#434's structure but pin the known intermediates
  ($q+q\to\bar{\rep3}$ diquark, diquark $+\,q\to$ proton, and conjugates); confirm the
  colored diquark is hosted in the bulk (low $r_U$, per \#416) and compare the emergent
  path A against the imposed path B.
  \item \textbf{Flavor: the symmetry-breaking isospin doublet.} The \#414 cure --- a broken
  $SU(2)$ that \emph{acts} ($p\leftrightarrow n$ ladder), a genuinely new construction, not
  a split of the existing register.
  \item \textbf{\#429's equivariant $d\ge3$ gap-fill.} Close the remaining
  coface-mirroring generalization so the dimension-generic apex is fully equivariant in
  $d\ge3$, completing the $3+1$\,D interior the \#418 spike needs.
  \item \textbf{Flux-tube creation energetics.} The energy cost of color-emergent creation
  and the flux-tube structure between the bipartite vertices --- the dynamical content the
  assembly experiments will probe.
\end{itemize}

\paragraph{Verdict.} The track delivered its geometric foundations and three of its four
sectors as merged, machine-verified results; it killed two hypotheses cleanly (flavor in
the existing register; the per-slice twist reading) and answered the binding question
(confinement is not fundamental in a connected bulk). The net is a sharper architecture ---
bipartite junctions, one global relaxation, color-emergent creation, charge in the temporal
sector and color in the register --- and a precise list of what to build next. As
throughout the program: we read what the geometry produces, we do not impose the answer,
and the falsifiable tests of \S\ref{sec:faithful} arbitrate.

\end{document}
